% !TeX TXS-program:compile = txs:///pdflatex

\documentclass[10pt]{article}
\usepackage[a4paper,margin=1cm]{geometry}
%\usepackage[bitstream-charter]{mathdesign}
\usepackage{xintexpr}
\usepackage{enumitem}
\usepackage{amsmath}
\usepackage{multicol}
\usepackage{minted-code}
\setlength\parindent{0pt}
\setlength\fboxsep{1pt}
\DeclareMathSymbol{;}{\mathbin}{operators}{'73}

\usepackage{tkzinterval}

\begin{document}

\pagestyle{empty}

\part*{Projet de macro \texttt{tkzinterval}}

\section{Utilisation}

\subsection{Macro}

\begin{codeblockA}[title=Macro et clés]
\tkzinterval[type,clés]{borne inf ; borne sup}
%le type peut-être FF (défaut) / OO / OF / FO
\end{codeblockA}

\subsection{Clés}

\begin{itemize}[noitemsep]
	\item \texttt{vsep} : espace vertical (en unité) entre la boîte et les crochets\hfill(\texttt{0.1em})
	\item \texttt{overlap} : effet 'superposition' (\texttt{lr-l-r-none}) pour les crochets ouverts\hfill(\texttt{lr})
	\item \texttt{hsep} : marge horizontale, en pourcentage de la largeur des crochets\hfill(\texttt{0.75})
	\item \texttt{space before} : espace avant la boîte\hfill(\texttt{\textbackslash,})
	\item \texttt{space after} : espace après la boîte\hfill(\texttt{\textbackslash,})
\end{itemize}

\subsection{Modification des paramètres}

\begin{codeblockA}[title=Gestion des paramètres de longueur]
\setlength\tkzintervhoriz{longueur}  %largeurs des crochets    (5pt)
\setlength\tkzintervwidth{longueur}  %épaisseur des crochets   (0.9pt)
\end{codeblockA}

\begin{codeblockA}[title=Gestion des paramètres globaux de la commande]
\setKVdefault[tkzinterv]{clé1 = nouvelle valeur par défaut,...}
\end{codeblockA}

\subsection{La clé \texttt{overlap}}

\begin{exampleM}{}
.\fbox{\tkzinterval[OO,overlap=lr]  {\dfrac{1}{2};\sqrt{7}}}.
.\fbox{\tkzinterval[OO,overlap=l]   {\dfrac{1}{2};\sqrt{7}}}.
.\fbox{\tkzinterval[OO,overlap=r]   {\dfrac{1}{2};\sqrt{7}}}.
.\fbox{\tkzinterval[OO,overlap=none]{\dfrac{1}{2};\sqrt{7}}}.
\end{exampleM}

\begin{exampleM}{}
\setKVdefault[tkzinterv]{space before={},space after={}}
.\fbox{\tkzinterval[OO,overlap=lr]  {\dfrac{1}{2};\sqrt{7}}}.
.\fbox{\tkzinterval[OO,overlap=l]   {\dfrac{1}{2};\sqrt{7}}}.
.\fbox{\tkzinterval[OO,overlap=r]   {\dfrac{1}{2};\sqrt{7}}}.
.\fbox{\tkzinterval[OO,overlap=none]{\dfrac{1}{2};\sqrt{7}}}.
\end{exampleM}

\begin{exampleM}{}
\setlength\tkzintervhoriz{7pt}\setKVdefault[tkzinterv]{space before={\,},space after={\,}}
.\fbox{\tkzinterval[OO,overlap=lr]  {\dfrac{1}{2};\sqrt{7}}}.
.\fbox{\tkzinterval[OO,overlap=l]   {\dfrac{1}{2};\sqrt{7}}}.
.\fbox{\tkzinterval[OO,overlap=r]   {\dfrac{1}{2};\sqrt{7}}}.
.\fbox{\tkzinterval[OO,overlap=none]{\dfrac{1}{2};\sqrt{7}}}.
\end{exampleM}

\pagebreak

\section{Influence de quelques paramètres}

\subsection{Valeur hsep=0.75}

\begin{codeblockA}[title=Gestion des paramètres]
\setlength\tkzintervhoriz{5pt}       %largeurs des traits H
\setlength\tkzintervwidth{0.9pt}     %épaisseur
\setKVdefault[tkzinterv]{hsep=0.75}  %espacement horizontal
\end{codeblockA}

\setlength\tkzintervhoriz{5pt}       %largeurs des traits H
\setlength\tkzintervwidth{0.9pt}     %épaisseur
\setKVdefault[tkzinterv]{hsep=0.75}  %espacement horizontal

\begin{exampleM}{}
On considère l'intervalle $\tkzinterval[OO,overlap=lr]{\dfrac{1}{2};\sqrt{7}}$ noté I.\par
On considère l'intervalle $\tkzinterval[OO,overlap=l]{\dfrac{1}{2};\sqrt{7}}$ noté I.\par
On considère l'intervalle $\tkzinterval[OO,overlap=r]{\dfrac{1}{2};\sqrt{7}}$ noté I.\par
On considère l'intervalle $\tkzinterval[OO,overlap=none]{\dfrac{1}{2};\sqrt{7}}$ noté I.
\end{exampleM}

\begin{exampleM}{}
On considère l'intervalle $\tkzinterval[FF]{\dfrac{1}{2};\sqrt{7}}$ noté I. \&\
On considère l'intervalle $\left[\dfrac{1}{2};\sqrt{7}\right]$ noté I.
\end{exampleM}

\begin{exampleM}{}
{\huge On considère l'intervalle $\mathcolor{red}{\tkzinterval[FF]{\dfrac{1}{2};\sqrt{7}}}$ noté I.}
\end{exampleM}

\begin{exampleM}{}
On considère l'intervalle $\tkzinterval[OF,overlap=l]{\dfrac{1}{2};\sqrt{7}}$ noté I.\par
On considère l'intervalle $\tkzinterval[OF,overlap=none]{\dfrac{1}{2};\sqrt{7}}$ noté I.
\end{exampleM}

\begin{exampleM}{}
On considère l'intervalle $\tkzinterval[FO,overlap=r]{\dfrac{1}{2};\sqrt{7}}$ noté I.\par
On considère l'intervalle $\tkzinterval[FO,overlap=none]{\dfrac{1}{2};\sqrt{7}}$ noté I.
\end{exampleM}

\pagebreak

\subsection{Valeur hsep=0.5}

\begin{codeblockA}[title=Gestion des paramètres]
\setlength\tkzintervhoriz{5pt}       %largeurs des traits H
\setlength\tkzintervwidth{0.9pt}     %épaisseur
\setKVdefault[tkzinterv]{hsep=0.5}   %espacement horizontal
\end{codeblockA}

\setlength\tkzintervhoriz{5pt}       %largeurs des traits H
\setlength\tkzintervwidth{0.9pt}     %épaisseur
\setKVdefault[tkzinterv]{hsep=0.5}   %espacement horizontal

\begin{exampleM}{}
On considère l'intervalle $\tkzinterval[OO,overlap=lr]{\dfrac{1}{2};\sqrt{7}}$ noté I.\par
On considère l'intervalle $\tkzinterval[OO,overlap=l]{\dfrac{1}{2};\sqrt{7}}$ noté I.\par
On considère l'intervalle $\tkzinterval[OO,overlap=r]{\dfrac{1}{2};\sqrt{7}}$ noté I.\par
On considère l'intervalle $\tkzinterval[OO,overlap=none]{\dfrac{1}{2};\sqrt{7}}$ noté I.
\end{exampleM}

\begin{exampleM}{}
On considère l'intervalle $\tkzinterval[FF]{\dfrac{1}{2};\sqrt{7}}$ noté I
\end{exampleM}

\begin{exampleM}{}
{\huge On considère l'intervalle $\tkzinterval[FF]{\dfrac{1}{2};\sqrt{7}}$ noté I}
\end{exampleM}

\begin{exampleM}{}
On considère l'intervalle $\tkzinterval[OF,overlap=l]{\dfrac{1}{2};\sqrt{7}}$ noté I.\par
On considère l'intervalle $\tkzinterval[OF,overlap=none]{\dfrac{1}{2};\sqrt{7}}$ noté I.
\end{exampleM}

\begin{exampleM}{}
On considère l'intervalle $\tkzinterval[FO,overlap=r]{\dfrac{1}{2};\sqrt{7}}$ noté I.\par
On considère l'intervalle $\tkzinterval[FO,overlap=none]{\dfrac{1}{2};\sqrt{7}}$ noté I.
\end{exampleM}

\pagebreak

\subsection{Valeurs hsep=1\&\ horiz=1mm \&\ width=0.25mm}

\begin{codeblockA}[title=Gestion des paramètres]
\setlength\tkzintervhoriz{1.75mm}    %largeurs des traits H
\setlength\tkzintervwidth{0.65mm}    %épaisseur
\setKVdefault[tkzinterv]{hsep=1}     %espacement horizontal
\end{codeblockA}

\setlength\tkzintervhoriz{1.75mm}
\setlength\tkzintervwidth{0.65mm}
\setKVdefault[tkzinterv]{hsep=1}

\begin{exampleM}{}
On considère l'intervalle $\tkzinterval[OO,overlap=lr]{\dfrac{1}{2};\sqrt{7}}$ noté I.\par
On considère l'intervalle $\tkzinterval[OO,overlap=l]{\dfrac{1}{2};\sqrt{7}}$ noté I.\par
On considère l'intervalle $\tkzinterval[OO,overlap=r]{\dfrac{1}{2};\sqrt{7}}$ noté I.\par
On considère l'intervalle $\tkzinterval[OO,overlap=none]{\dfrac{1}{2};\sqrt{7}}$ noté I.
\end{exampleM}

\begin{exampleM}{}
On considère l'intervalle $\tkzinterval[FF]{\dfrac{1}{2};\sqrt{7}}$ noté I
\end{exampleM}

\begin{exampleM}{}
{\huge On considère l'intervalle $\tkzinterval[FF]{\dfrac{1}{2};\sqrt{7}}$ noté I}
\end{exampleM}

\begin{exampleM}{}
On considère l'intervalle $\tkzinterval[OF,overlap=l]{\dfrac{1}{2};\sqrt{7}}$ noté I.\par
On considère l'intervalle $\tkzinterval[OF,overlap=none]{\dfrac{1}{2};\sqrt{7}}$ noté I.
\end{exampleM}

\begin{exampleM}{}
On considère l'intervalle $\tkzinterval[FO,overlap=r]{\dfrac{1}{2};\sqrt{7}}$ noté I.\par
On considère l'intervalle $\tkzinterval[FO,overlap=none]{\dfrac{1}{2};\sqrt{7}}$ noté I.
\end{exampleM}

%------------------

\pagebreak

\section{Utilisation 'poussée' ?}

\subsection{Avec des réunions / intersections}

\begin{codeblockA}[title=Gestion des paramètres]
\setlength\tkzintervhoriz{5pt}       %largeurs des traits H
\setlength\tkzintervwidth{0.8pt}     %épaisseur
\setKVdefault[tkzinterv]{hsep=0.75}  %espacement horizontal
\end{codeblockA}

\setlength\tkzintervhoriz{5pt}
\setlength\tkzintervwidth{0.8pt}
\setKVdefault[tkzinterv]{hsep=0.75}  %espacement horizontal

\begin{exampleM}{}
$x \in \tkzinterval[FF]{4;\dfrac{100}{7}} \cup \tkzinterval[OO]{\dfrac{1}{2};\sqrt{7}}$.\hfill
$x \in \tkzinterval[OO]{4;\dfrac{100}{7}} \cup \tkzinterval[OO]{\dfrac{1}{2};\sqrt{7}}$.\hfill\null

\hfill$x \in \tkzinterval[FO]{4;\dfrac{100}{7}} \cup \tkzinterval[OO]{\dfrac{1}{2};\sqrt{7}}$.\hfill
$x \in \tkzinterval[OF]{4;\dfrac{100}{7}} \cup \tkzinterval[OO]{\dfrac{1}{2};\sqrt{7}}$.
\end{exampleM}

\begin{exampleM}{}
$x \in \tkzinterval[FF]{4;\dfrac{100}{7}} \cup \tkzinterval[FF]{\dfrac{1}{2};\sqrt{7}}$.\hfill
$x \in \tkzinterval[OO]{4;\dfrac{100}{7}} \cup \tkzinterval[FF]{\dfrac{1}{2};\sqrt{7}}$.\hfill\null

\hfill$x \in \tkzinterval[FO]{4;\dfrac{100}{7}} \cup \tkzinterval[FF]{\dfrac{1}{2};\sqrt{7}}$.\hfill
$x \in \tkzinterval[OF]{4;\dfrac{100}{7}} \cup \tkzinterval[FF]{\dfrac{1}{2};\sqrt{7}}$.
\end{exampleM}

\begin{exampleM}{}
$x \in \tkzinterval[FF]{4;\dfrac{100}{7}} \cup \tkzinterval[OF]{\dfrac{1}{2};\sqrt{7}}$.\hfill
$x \in \tkzinterval[OO]{4;\dfrac{100}{7}} \cup \tkzinterval[OF]{\dfrac{1}{2};\sqrt{7}}$.\hfill\null

\hfill$x \in \tkzinterval[FO]{4;\dfrac{100}{7}} \cup \tkzinterval[OF]{\dfrac{1}{2};\sqrt{7}}$.\hfill
$x \in \tkzinterval[OF]{4;\dfrac{100}{7}} \cup \tkzinterval[OF]{\dfrac{1}{2};\sqrt{7}}$.
\end{exampleM}

\begin{exampleM}{}
$x \in \tkzinterval[FF]{4;\dfrac{100}{7}} \cup \tkzinterval[FO]{\dfrac{1}{2};\sqrt{7}}$.\hfill
$x \in \tkzinterval[OO]{4;\dfrac{100}{7}} \cup \tkzinterval[FO]{\dfrac{1}{2};\sqrt{7}}$.\hfill\null

\hfill$x \in \tkzinterval[FO]{4;\dfrac{100}{7}} \cup \tkzinterval[FO]{\dfrac{1}{2};\sqrt{7}}$.\hfill
$x \in \tkzinterval[OF]{4;\dfrac{100}{7}} \cup \tkzinterval[FO]{\dfrac{1}{2};\sqrt{7}}$.
\end{exampleM}

\end{document}